% How accurate LIGHT v0_1 is: its accuracy, its system map and its reaction
% rates (Appendix A). The v0 overview (accuracy_overview.tex) describes
% v0. No number here is typed by hand. They come from five generated files:
%   system_map_v0_1_numbers.tex    make_system_map_v0_1.py
%   overview_v0_1_rates.tex, overview_v0_1_fast.tex, overview_v0_1_slow.tex
%                                  make_overview_v0_1.py
%   energy_ladder_v0_1.tex         make_energy_ladder.py
% Build from LIGHT/docs:
%   python3 make_system_map_v0_1.py && python3 make_overview_v0_1.py
%   python3 make_energy_ladder.py
%   latexmk -pdf system_map_v0_1.tex && latexmk -pdf accuracy_overview_v0_1.tex
\documentclass[11pt]{article}
\usepackage[a4paper,margin=2.2cm]{geometry}
\usepackage[T1]{fontenc}
\usepackage[scaled=0.95]{helvet}
\renewcommand{\familydefault}{\sfdefault}
\usepackage{amsmath}
\usepackage{sansmath}
\usepackage{booktabs,tabularx,array,longtable}
\usepackage{xcolor}
\usepackage{pdfpages}
\usepackage{float}
\usepackage{enumitem}
\usepackage{needspace}
\usepackage{tikz}
\usetikzlibrary{arrows.meta,decorations.pathmorphing}
\setlist{nosep, leftmargin=1.2em}
\setlength{\parindent}{0pt}
\setlength{\parskip}{5pt}
\definecolor{ink2}{HTML}{52514E}

% LIGHT v0_1 map numbers. GENERATED by make_system_map_v0_1.py
% from the model (0.1.1). Do not edit. Rerun the script instead.
\providecommand\vmdef[2]{\expandafter\def\csname vm@#1\endcsname{#2}}
\vmdef{e.P}{100.0}
\vmdef{e.pen}{31.6}
\vmdef{e.G1}{96.6}
\vmdef{e.C1}{14.2}
\vmdef{e.C2}{0.5}
\vmdef{e.F2}{2.8}
\vmdef{e.G6}{9.4}
\vmdef{e.dU}{0.0}
\vmdef{e.cap}{8.1}
\vmdef{C1.kr}{12.1}
\vmdef{C1.heat}{2.1}
\vmdef{F2e.F}{1.5}
\vmdef{F2e.heat}{2.1}
\vmdef{F2e.anion}{-0.7}
\vmdef{krs.in}{160.6}
\vmdef{krs.harp}{51.7}
\vmdef{krs.G2}{107.6}
\vmdef{krs.left}{1.4}
\vmdef{brs.harp}{32.2}
\vmdef{brs.G2}{67.0}
\vmdef{harp.B}{17.0}
\vmdef{harp.C}{10.4}
\vmdef{harp.F}{4.6}
\vmdef{harp.heat}{19.6}
\vmdef{G2.ion}{75.9}
\vmdef{C2.ion}{0.5}
\vmdef{K15.heat}{6.2}
\vmdef{R11p.in}{67.2}
\vmdef{R11p.kr}{51.9}
\vmdef{R11p.heat}{15.3}
\vmdef{ii.in}{3.7}
\vmdef{ii.anion}{-0.7}
\vmdef{ii.B}{1.6}
\vmdef{ii.heat}{1.4}
\vmdef{ions.left}{-0.6}
\vmdef{anion.left}{0.0}
\vmdef{mix.net}{3.1}
\vmdef{mix.netC}{3.0}
\vmdef{mix.heat}{0.0}
\vmdef{mix.BC}{10.2}
\vmdef{mix.CB}{13.3}
\vmdef{B.G4}{5.5}
\vmdef{B.G4A}{0.2}
\vmdef{B.G5}{14.1}
\vmdef{B.Q}{1.9}
\vmdef{B.QF2}{0.9}
\vmdef{B.QKR}{0.0}
\vmdef{B.QE}{1.0}
\vmdef{B.left}{0.0}
\vmdef{C.G4C}{1.0}
\vmdef{C.G5C}{3.7}
\vmdef{C.Q}{2.7}
\vmdef{C.QF2}{0.7}
\vmdef{C.QKR}{1.0}
\vmdef{C.QE}{1.0}
\vmdef{C.left}{0.0}
\vmdef{G4.hv}{4.7}
\vmdef{G4.F}{0.8}
\vmdef{G4.heat}{0.1}
\vmdef{G4A.hv}{0.2}
\vmdef{G4A.F}{0.0}
\vmdef{G4A.heat}{0.0}
\vmdef{G4C.hv}{0.7}
\vmdef{G4C.F}{0.1}
\vmdef{G4C.heat}{0.1}
\vmdef{QB.F}{0.5}
\vmdef{QB.heat}{1.4}
\vmdef{QC.F}{0.6}
\vmdef{QC.heat}{2.1}
\vmdef{tri.k2f}{16.0}
\vmdef{tri.heat}{1.8}
\vmdef{k2f.C3}{8.2}
\vmdef{k2f.Q2F}{7.8}
\vmdef{k2f.left}{0.0}
\vmdef{brk.C3}{51.1}
\vmdef{brk.Q2F}{48.9}
\vmdef{C3.hv}{4.6}
\vmdef{C3.F}{1.2}
\vmdef{C3.heat}{2.4}
\vmdef{Q2F.F}{1.2}
\vmdef{Q2F.heat}{6.6}
\vmdef{FF.in}{0.1}
\vmdef{FF.heat}{0.1}
\vmdef{out.248}{4.7}
\vmdef{out.272}{0.2}
\vmdef{out.275}{0.7}
\vmdef{out.420}{4.6}
\vmdef{out.heat}{69.2}
\vmdef{out.elastic}{9.4}
\vmdef{out.ions}{-0.6}
\vmdef{out.F}{10.4}
\vmdef{out.other}{1.4}
\vmdef{G4x.F}{0.8}
\vmdef{G4x.heat}{0.1}
\vmdef{version}{0.1.1}
\vmdef{p}{1}
\vmdef{xF2}{0.5}
\vmdef{Tg0}{300}
\vmdef{P}{300}
\vmdef{pulse}{1}
\vmdef{tend}{2}
\vmdef{ne}{10^{15}}
\vmdef{neF2}{1.95{\times}10^{14}}
\vmdef{nePeak}{3.66{\times}10^{15}}
\vmdef{N}{2.45{\times}10^{19}}
\vmdef{Edep}{0.30}
\vmdef{Te}{4.05}
\vmdef{TeMid}{3.25}
\vmdef{rTail}{0.18}
\vmdef{TtMid}{0.59}
\vmdef{TgEnd}{762}
\vmdef{pEnd}{2.55}
\vmdef{F2burn}{613}
\vmdef{F2left}{0.09}
\vmdef{krsFormed}{3.03{\times}10^{17}}
\vmdef{krfFormed}{9.30{\times}10^{16}}
\vmdef{k2fFormed}{5.69{\times}10^{16}}
\vmdef{krf.eff}{24.8}
\vmdef{frm.ii}{5.6}
\vmdef{brf.G4}{19.1}
\vmdef{brf.G4A}{0.7}
\vmdef{brf.G4C}{3.3}
\vmdef{brf.tri}{61.2}
\vmdef{brf.e}{6.8}
\vmdef{brf.F2}{5.6}
\vmdef{brf.Kr}{3.4}
\vmdef{brf.QB}{6.4}
\vmdef{brf.QC}{9.4}
\vmdef{ion.R11p}{95.0}
\vmdef{ion.ii}{5.0}
\vmdef{k.harp}{8.47{\times}10^{-10}}
\vmdef{pc.harpY}{90}
\vmdef{pc.harpB}{56}
\vmdef{pc.harpC}{34}
\vmdef{pc.harpX}{10}
\vmdef{pc.fC}{38}
\vmdef{k.G2}{5.0{\times}10^{-10}}
\vmdef{tauB}{6.99}
\vmdef{pc.BA}{3.5}
\vmdef{tauC}{50}
\vmdef{tauC3}{220}
\vmdef{gapBC}{120}
\vmdef{k.BC}{1.05{\times}10^{-11}}
\vmdef{Keq}{1.10}
\vmdef{k.G5}{1.9{\times}10^{-30}}
\vmdef{k.G5C}{3.79{\times}10^{-31}}
\vmdef{tpow.G5}{-3.32}
\vmdef{k.Q2F}{1.2{\times}10^{-10}}
\vmdef{k.K15}{2.4{\times}10^{-31}}
\vmdef{tpow.K15}{-0.75}
\vmdef{k.R11p}{1.6{\times}10^{-6}}
\vmdef{pow.R11p}{0.55}
\vmdef{cap.R11p}{0.95}
\vmdef{k.FF}{2.8{\times}10^{-34}}
\vmdef{n3.ii}{-3}
\vmdef{E.anion}{2.6}
\vmdef{hv.G4C}{4.51}
\vmdef{hv.G4}{5.00}
\vmdef{hv.G4A}{4.56}
\vmdef{hv.C3}{2.95}
\vmdef{E.B}{5.856}
\vmdef{notes.krs}{110.8}
\vmdef{notes.krp}{76.5}
\vmdef{notes.emit}{10.2}
\vmdef{om.kr2}{9.2}
\vmdef{om.electron}{4.5}
\vmdef{om.all}{6.6}
\vmdef{om.basic}{4.7}
\vmdef{om.kr2noF}{4.7}
\vmdef{om.v0}{16.5}
\vmdef{om.trimer}{2.3}

% LIGHT v0_1 reaction-rate tables. GENERATED by make_overview_v0_1.py
% from LIGHT v0_1 (0.1.1). Do not edit. Rerun the script instead.
\providecommand\videf[2]{\expandafter\def\csname vi@#1\endcsname{#2}}
\videf{rates.teMid}{3.25}
\videf{rates.rMid}{0.18}
\videf{rates.nBasic}{37}
\videf{rates.nGroups}{8}
\videf{rates.groups}{kr2, electron}
\videf{rates.v0groups}{quench, kr2, ions, electron, bc, f2e}
\def\virBasic{G1 & \el{} + Kr \ra \el{} + \Krs{} & $\langle\sigma v\rangle$, Table A.3 & $T_e$ & Biagi v7.1, the whole 5s manifold at 9.915 eV\\
C1a & \el{} + Kr \ra \el{} + \Krs{} & $\langle\sigma v\rangle$, Table A.3 & $T_e$ & Biagi v7.1 5p group at 11.30 eV, prompt cascade to 5s\\
C1b & \el{} + Kr \ra \el{} + \Krs{} & $\langle\sigma v\rangle$, Table A.3 & $T_e$ & Biagi v7.1 d+p group at 11.998 eV, as C1a\\
C1c & \el{} + Kr \ra \el{} + \Krs{} & $\langle\sigma v\rangle$, Table A.3 & $T_e$ & Biagi v7.1 higher group at 12.75 eV, as C1a\\
C2 & \el{} + Kr \ra 2\el{} + \Krp{} & $\langle\sigma v\rangle$, Table A.3 & $T_e$ & Biagi v7.1 ionization at 13.996 eV\\
G6 & \el{} + Kr, elastic & energy loss $(2m/M)\langle\sigma_m v\varepsilon\rangle$ & $T_e$, $T_g$ & Biagi v7.1 momentum transfer\\
ATT & \el{} + \Ft{} \ra F\ensuremath{^-} + F & $\langle\sigma v\rangle$, Table A.3 & $T_e$ & SIGLO F$_2$ set\\
DIS\_a3 & \el{} + \Ft{} \ra \el{} + 2F & $\langle\sigma v\rangle$, Table A.3 & $T_e$ & SIGLO {\unsansmath$a\,^3\Pi_u$}, repulsive, so excitation dissociates\\
DIS\_A1 & \el{} + \Ft{} \ra \el{} + 2F & $\langle\sigma v\rangle$, Table A.3 & $T_e$ & SIGLO {\unsansmath$A\,^1\Pi_u$}, as DIS\_a3\\
VIB1 & \el{} + \Ft{} \ra \el{} + \Ft{} & $\langle\sigma v\rangle$, Table A.3 & $T_e$ & SIGLO, v = 1\\
VIB2 & \el{} + \Ft{} \ra \el{} + \Ft{} & $\langle\sigma v\rangle$, Table A.3 & $T_e$ & SIGLO, v = 2\\
VIB3 & \el{} + \Ft{} \ra \el{} + \Ft{} & $\langle\sigma v\rangle$, Table A.3 & $T_e$ & SIGLO, v = 3\\
VIB4 & \el{} + \Ft{} \ra \el{} + \Ft{} & $\langle\sigma v\rangle$, Table A.3 & $T_e$ & SIGLO, v = 4\\
G2 & 2\Krs{} \ra \Krp{} + Kr + \el{} & $5.0{\times}10^{-10}$ cm$^3$s$^{-1}$ &  & Kannari 1985 Table I, Penning block\\
G3 & \Krs{} + \Ft{} \ra KrF(B) + F & $4.73{\times}10^{-10}$ cm$^3$s$^{-1}$ &  & rate from Velazco, Kolts and Setser 1976 and Ueda 1984, KrF$^*$ yield 0.90 from Velazco, Kolts and Setser 1976, B share from Yu and Setser 1990\\
G3C & \Krs{} + \Ft{} \ra KrF(C) + F & $2.9{\times}10^{-10}$ cm$^3$s$^{-1}$ &  & as G3, with Yu and Setser's $k_B/k_C$ of 1.6\\
G3X & \Krs{} + \Ft{} \ra Kr + 2F & $8.47{\times}10^{-11}$ cm$^3$s$^{-1}$ &  & as G3, the KrF$^*$-free rest, assigned to Kr + 2F\\
K15 & \Krp{} + 2Kr \ra \Krtp{} + Kr & $2.4{\times}10^{-31}$ cm$^6$s$^{-1}$ & $(T_g/300)^{-0.75}$ & the mean of three measurements, 2.4$\times10^{-31}$. $T_g$ law from Perez-Rios and Greene 2018\\
R11p & \el{} + \Krtp{} \ra \Krs{} + Kr & $1.6{\times}10^{-6}$ $(T_e/300\,\mathrm{K})^{-0.55}$ cm$^3$s$^{-1}$ & $T_e^{-0.55}$ & Shiu and Biondi 1977, 1.6$\times10^{-6}$ $(300\,\mathrm{K}/T_e)^{0.55}$, $\pm$ 10\,\%. Mikus 1978 within 8\,\%\\
K26 & \Krp{} + F\ensuremath{^-} \ra KrF(B) & $2.95{\times}10^{-6}$ cm$^3$s$^{-1}$ at 1\,atm, 300\,K, unscreened & $N$, screening, $(T_g/300)^{-3}$ on the 3-body part & Kannari 1985, Flannery turnover with Debye screening, and $(T_g/300)^{-3}$ on its three-body part from Thomson theory. Measured with F$^-$ only below 300 Torr, Jacob et al. 1978\\
K108 & \Krtp{} + F\ensuremath{^-} \ra KrF(B) + Kr & $2.95{\times}10^{-6}$ cm$^3$s$^{-1}$ at 1\,atm, 300\,K, unscreened & $N$, screening, $(T_g/300)^{-3}$ on the 3-body part & as K26\\
G4 & KrF(B) \ra Kr + F + \hnu & $1.38{\times}10^{8}$ s$^{-1}$ &  & Kannari 1985 Table I, the 96.6\,\% B-X share of $\tau_B$ = 6.99 ns. Eden 1978 measured 6.8 $\pm$ 0.2 ns\\
G4A & KrF(B) \ra Kr + F + \hnu & $4.84{\times}10^{6}$ s$^{-1}$ &  & $A_{B\text{-}A}/A_{B\text{-}X}$ = 0.035, Dunning and Hay's calculation via Yu and Setser 1990. 272 nm is a theoretical position\\
G4C & KrF(C) \ra Kr + F + \hnu & $2.0{\times}10^{7}$ s$^{-1}$ &  & ML joint fit. Hakuta 1990 measured 60 $\pm$ 10 ns. Band at 275 nm, printed\\
BC & KrF(B) + Kr \ra KrF(C) + Kr & $1.05{\times}10^{-11}$ cm$^3$s$^{-1}$ &  & ML joint fit. Hakuta 1990's 1.2$\times10^{-10}$ is held out\\
CB & KrF(C) + Kr \ra KrF(B) + Kr & $9.55{\times}10^{-12}$ cm$^3$s$^{-1}$ & $1/K_\mathrm{eq}(T_g)$ & detailed balance of BC, with the joint fit's $K_\mathrm{eq}$ and Yu and Setser's measured 120 $\pm$ 50 cm$^{-1}$ gap\\
G5 & KrF(B) + 2Kr \ra \KtF{} + Kr & $1.9{\times}10^{-30}$ cm$^6$s$^{-1}$ & $(T_g/300)^{-3.32}$ & ML joint fit. $T_g$ law from Shui's calculation for 300-800 K, measured only to 425 K\\
G5C & KrF(C) + 2Kr \ra \KtF{} + Kr & $3.79{\times}10^{-31}$ cm$^6$s$^{-1}$ & $(T_g/300)^{-3.32}$ & ML joint fit\\
QE & KrF(B) + \el{} \ra Kr + F + \el{} & $6.6{\times}10^{-8}$ cm$^3$s$^{-1}$ &  & Trainor and Jacob's measured product, re-read with this network's lifetime\\
QEC & KrF(C) + \el{} \ra Kr + F + \el{} & $6.6{\times}10^{-8}$ cm$^3$s$^{-1}$ &  & as QE\\
QF2 & KrF(B) + \Ft{} \ra Kr + 3F & $4.75{\times}10^{-10}$ cm$^3$s$^{-1}$ &  & ML joint fit. Xu 1993 printed 3.0 $\pm$ 0.3$\times10^{-10}$ with fast mixing\\
QF2C & KrF(C) + \Ft{} \ra Kr + 3F & $3.04{\times}10^{-10}$ cm$^3$s$^{-1}$ &  & ML joint fit\\
QKR & KrF(B) + Kr \ra 2Kr + F & $1.0{\times}10^{-14}$ cm$^3$s$^{-1}$ &  & ML joint fit, at its lower bound\\
QKRC & KrF(C) + Kr \ra 2Kr + F & $8.37{\times}10^{-13}$ cm$^3$s$^{-1}$ &  & ML joint fit\\
C3 & \KtF{} \ra 2Kr + F + \hnu & $4.5{\times}10^{6}$ s$^{-1}$ &  & Xu, Gadomski and Setser 1993, 220 $\pm$ 20 ns. Band at 420 nm, printed\\
Q2F & \KtF{} + \Ft{} \ra 2Kr + F + \Ft{} & $1.2{\times}10^{-10}$ cm$^3$s$^{-1}$ &  & Xu 1993 Table I, 1.2 $\pm$ 0.2$\times10^{-10}$\\
FF & 2F + Kr \ra \Ft{} + Kr & $2.8{\times}10^{-34}$ cm$^6$s$^{-1}$ &  & Cohen and Bott 1976 evaluation, 2.8$\times10^{-34}$, every rare gas given the Ar value. Held at its 300 K value\\}
\def\virGroups{AK8\_Kr & \Krs{} + 2Kr \ra \Krts{} + Kr & $1.7{\times}10^{-32}$ cm$^6$s$^{-1}$ & $(T_g/300)^{-0.64}$ & Ueda 1984 Table V, 1.7 $\pm$ 0.2$\times10^{-32}$. $T_g$ law from Wieme and Lenaerts 1981\\
KR2\_RAD & \Krts{} \ra 2Kr + \hnu & $3.7{\times}10^{6}$ s$^{-1}$ &  & Morikawa 1989, 270 $\pm$ 5 ns, and five other measurements\\
NK7 & \Krts{} + \Ft{} \ra KrF(B) + Kr + F & $1.05{\times}10^{-10}$ cm$^3$s$^{-1}$ &  & half of Quigley and Hughes 1978, 2.1 $\pm$ 0.2$\times10^{-10}$. The even split is Kannari's\\
K2F\_BR & \Krts{} + \Ft{} \ra \KtF{} + F & $1.05{\times}10^{-10}$ cm$^3$s$^{-1}$ &  & the other half of Quigley and Hughes 1978\\
K2F\_F & \Krts{} + F \ra KrF(B) + Kr & $3.0{\times}10^{-10}$ cm$^3$s$^{-1}$ &  & Kannari 1985, an estimate. No measurement exists\\
PEN\_KR2 & 2\Krts{} \ra \Krtp{} + 2Kr + \el{} & $5.0{\times}10^{-10}$ cm$^3$s$^{-1}$ &  & Kannari 1985 Table I, from an unpublished 1973 SRI report that gives Ar$_2^*$ the same value. No measurement held\\
KSTEP & \el{} + \Krs{} \ra \Krp{} + 2\el{} & $1.56{\times}10^{-7}\,T_e^{0.71}\times$ $e^{-4.09/T_e}$ cm$^3$s$^{-1}$ & $T_e$ & Chen R9 template, an estimate\\
SUPER & \el{} + \Krs{} \ra \el{} + Kr & from G1 by detailed balance & $T_e$ & detailed balance of G1, statistical weight ratio 1/12\\}
\def\virVzero{G3 & \Krs{} + \Ft{} \ra KrF(B) + F & $8.1{\times}10^{-10}$ cm$^3$s$^{-1}$ &  & Kannari 1985 Table I, all to B. The report it credits gives 7.2$\times10^{-10}$\\
K15 & \Krp{} + 2Kr \ra \Krtp{} + Kr & $2.5{\times}10^{-31}$ cm$^6$s$^{-1}$ & $(T_g/300)^{-0.75}$ & Kannari 1985, Boichenko 2000\\
R11p & \el{} + \Krtp{} \ra \Krs{} + Kr & $2.24{\times}10^{-7}\,T_e^{-1/2}$ cm$^3$s$^{-1}$ & $T_e$ & Kannari 1985 Table I, 1.39$\times10^{-6}$ at 300 K on a $T_e^{-1/2}$ template\\
K26 & \Krp{} + F\ensuremath{^-} \ra KrF(B) & $2.95{\times}10^{-6}$ cm$^3$s$^{-1}$ at 1\,atm, 300\,K, unscreened & $N$, screening & Kannari 1985, Flannery turnover with Debye screening, no $T_g$ law\\
K108 & \Krtp{} + F\ensuremath{^-} \ra KrF(B) + Kr & $2.95{\times}10^{-6}$ cm$^3$s$^{-1}$ at 1\,atm, 300\,K, unscreened & $N$, screening & as K26\\
G4 & KrF(B) \ra Kr + F + \hnu & $1.43{\times}10^{8}$ s$^{-1}$ &  & Kannari 1985 Table I, all of $\tau_B$ as B-X\\
G5 & KrF(B) + 2Kr \ra \KtF{} + Kr & $6.0{\times}10^{-31}$ cm$^6$s$^{-1}$ & $(T_g/300)^{-3.32}$ & Xu, Gadomski and Setser 1993 Table II, 6.0 $\pm$ 2.0$\times10^{-31}$\\
G5\_C & KrF(C) + 2Kr \ra \KtF{} + Kr & $6.0{\times}10^{-31}$ cm$^6$s$^{-1}$ & $(T_g/300)^{-3.32}$ & as G5, applied to KrF(C)\\
AK8\_Kr & \Krs{} + 2Kr \ra \Krts{} + Kr & $2.7{\times}10^{-32}$ cm$^6$s$^{-1}$ &  & REBUILD's estimate, Kannari's Ar-partner rate given to Kr\\
KR2\_RAD & \Krts{} \ra 2Kr + \hnu & $8.0{\times}10^{7}$ s$^{-1}$ &  & Kannari 1985 Table I, 8.0$\times10^{7}$ per second\\
NK7 & \Krts{} + \Ft{} \ra KrF(B) + Kr + F & $3.0{\times}10^{-10}$ cm$^3$s$^{-1}$ &  & Kannari 1985 Table I\\
K2F\_BR & \Krts{} + \Ft{} \ra \KtF{} + F & $3.0{\times}10^{-10}$ cm$^3$s$^{-1}$ &  & Kannari 1985 Table I\\
K33 & KrF(B) + Kr \ra 2Kr + F & $4.0{\times}10^{-12}$ cm$^3$s$^{-1}$ &  & Xu, Gadomski and Setser 1993, 4.0 $\pm$ 0.5$\times10^{-12}$\\
K33\_C & KrF(C) + Kr \ra 2Kr + F & $4.0{\times}10^{-12}$ cm$^3$s$^{-1}$ &  & as K33, applied to KrF(C)\\
K32 & KrF(B) + \Ft{} \ra Kr + 3F & $3.0{\times}10^{-10}$ cm$^3$s$^{-1}$ &  & Xu, Gadomski and Setser 1993, 3.0 $\pm$ 0.3$\times10^{-10}$\\
K32\_C & KrF(C) + \Ft{} \ra Kr + 3F & $3.0{\times}10^{-10}$ cm$^3$s$^{-1}$ &  & as K32, applied to KrF(C)\\
K67 & KrF(B) + \el{} \ra Kr + F + \el{} & $4.5{\times}10^{-8}$ cm$^3$s$^{-1}$ &  & Kushner, Hanson and Schneider 1989\\
K67\_C & KrF(C) + \el{} \ra Kr + F + \el{} & $4.5{\times}10^{-8}$ cm$^3$s$^{-1}$ &  & as K67, applied to KrF(C)\\
K125 & \KtF{} + \Ft{} \ra 2Kr + F + \Ft{} & $1.5{\times}10^{-10}$ cm$^3$s$^{-1}$ &  & Kannari 1985 Table I\\
K\_2Fe & \KtF{} + \el{} \ra 2Kr + F + \el{} & $1.0{\times}10^{-7}$ cm$^3$s$^{-1}$ &  & Kannari 1985, estimated there\\
BC\_Kr & KrF(B) + Kr \ra KrF(C) + Kr & $5.0{\times}10^{-10}$ cm$^3$s$^{-1}$ &  & Kannari's B-C mixing with a Kr partner, an estimate\\
CB\_Kr & KrF(C) + Kr \ra KrF(B) + Kr & $3.41{\times}10^{-10}$ cm$^3$s$^{-1}$ &  & detailed balance of BC\_Kr at 300 K\\
URAD\_C & KrF(C) \ra Kr + F + \hnu & $1.33{\times}10^{7}$ s$^{-1}$ &  & Kannari 1985 Table I, 1.33$\times10^{7}$ per second\\}
\def\virXsec{G1 & Kr, 5s manifold, 9.915\,eV & $1.24{\times}10^{-17}$ & $3.44{\times}10^{-13}$ & $1.11{\times}10^{-11}$ & $6.41{\times}10^{-11}$ & $3.78{\times}10^{-10}$ & 0.044\\
C1a & Kr, 5p group, 11.3\,eV & $5.1{\times}10^{-19}$ & $5.93{\times}10^{-14}$ & $2.97{\times}10^{-12}$ & $2.11{\times}10^{-11}$ & $1.48{\times}10^{-10}$ & $7.03{\times}10^{-3}$\\
C1b & Kr, d+p group, 11.998\,eV & $2.5{\times}10^{-19}$ & $5.84{\times}10^{-14}$ & $3.77{\times}10^{-12}$ & $3.08{\times}10^{-11}$ & $2.53{\times}10^{-10}$ & $2.41{\times}10^{-3}$\\
C1c & Kr, higher group, 12.75\,eV & $9.15{\times}10^{-21}$ & $7.73{\times}10^{-15}$ & $8.8{\times}10^{-13}$ & $1.02{\times}10^{-11}$ & $1.29{\times}10^{-10}$ & $2.27{\times}10^{-4}$\\
C2 & Kr, ionization, 13.996\,eV & $8.06{\times}10^{-21}$ & $1.57{\times}10^{-14}$ & $2.28{\times}10^{-12}$ & $2.94{\times}10^{-11}$ & $4.14{\times}10^{-10}$ & $7.99{\times}10^{-5}$\\
ATT & F$_2$, attachment & $6.58{\times}10^{-9}$ & $3.51{\times}10^{-9}$ & $2.25{\times}10^{-9}$ & $1.61{\times}10^{-9}$ & $9.75{\times}10^{-10}$ & 1\\
DIS\_a3 & F$_2$, {\unsansmath$a\,^3\Pi_u$}, 3.16\,eV & $3.01{\times}10^{-12}$ & $1.13{\times}10^{-10}$ & $3.95{\times}10^{-10}$ & $7.37{\times}10^{-10}$ & $1.35{\times}10^{-9}$ & 1\\
DIS\_A1 & F$_2$, {\unsansmath$A\,^1\Pi_u$}, 4.34\,eV & $2.06{\times}10^{-13}$ & $1.93{\times}10^{-11}$ & $8.88{\times}10^{-11}$ & $1.9{\times}10^{-10}$ & $4.02{\times}10^{-10}$ & 1\\
VIB1 & F$_2$, $v=1$, 0.111\,eV & $1.77{\times}10^{-9}$ & $1.22{\times}10^{-9}$ & $8.72{\times}10^{-10}$ & $6.61{\times}10^{-10}$ & $4.27{\times}10^{-10}$ & 1\\
VIB2 & F$_2$, $v=2$, 0.219\,eV & $2.01{\times}10^{-9}$ & $1.49{\times}10^{-9}$ & $1.09{\times}10^{-9}$ & $8.37{\times}10^{-10}$ & $5.47{\times}10^{-10}$ & 1\\
VIB3 & F$_2$, $v=3$, 0.324\,eV & $1.21{\times}10^{-9}$ & $9.75{\times}10^{-10}$ & $7.35{\times}10^{-10}$ & $5.71{\times}10^{-10}$ & $3.78{\times}10^{-10}$ & 1\\
VIB4 & F$_2$, $v=4$, 0.421\,eV & $3.83{\times}10^{-10}$ & $3.27{\times}10^{-10}$ & $2.51{\times}10^{-10}$ & $1.96{\times}10^{-10}$ & $1.31{\times}10^{-10}$ & 1\\
\midrule
ATT & energy an attached \el{} takes, eV & 0.37 & 0.46 & 0.53 & 0.59 & 0.72 & \\
 & $\tfrac32 T_e$, for comparison & 0.75 & 1.50 & 2.25 & 3.00 & 4.50 & \\}
\def\virTe{0.5 & 1 & 1.5 & 2 & 3}
\def\virSpecies{Kr & 0 & 1.5 & ground state\\
\Ft{} & 0 & 2.5 & ground state\\
F & 0.8 & 1.5 & half the F$_2$ bond, $D_0$ = 1.60 eV\\
\Krs{} & 9.915 & 1.5 & Biagi's 5s threshold, the 5s[3/2]2 level\\
\Krp{} & 13.996 & 1.5 & Biagi's ionization threshold\\
KrF(B) & 5.856 & 2.5 & $D_e(B)$ = 5.54 eV below the ion pair, Tellinghuisen 1976\\
KrF(C) & 5.84112 & 2.5 & 120 cm$^{-1}$ below B, Yu and Setser 1990\\
\KtF{} & 5.276 & 3 & 0.58 eV below KrF$^*$ + Kr, calculated by Huestis et al. 1984\\
\Krts{} & 9.2 & 2.5 & the 147-nm excimer, an estimated well\\
\Krtp{} & 12.846 & 2.5 & IP(Kr) less $D_0$(Kr$_2^+$) = 1.15 eV\\
F\ensuremath{^-} & \ensuremath{-}2.6 & 1.5 & $E_\mathrm{form}$(F) less the electron affinity, 3.40 eV\\}

% LIGHT v0_1 accuracy numbers (fast). GENERATED by make_overview_v0_1.py
% from LIGHT v0_1 (0.1.1) and ML/. Do not edit. Rerun the script instead.
\providecommand\videf[2]{\expandafter\def\csname vi@#1\endcsname{#2}}
\videf{xu.50}{0.91}
\videf{xu.50.v0}{2.89}
\videf{xu.100}{0.97}
\videf{xu.100.v0}{2.41}
\videf{xu.50.meas}{$5.0{\times}10^{7}$}
\videf{xu.100.meas}{$6.2{\times}10^{7}$}
\videf{eq5.dev}{$1.6{\times}10^{-4}$}
\videf{gap}{120}
\videf{closure.max}{$10^{-12}$}
\videf{closure.runs}{9}
\videf{hot.p}{0.1}
\videf{hot.P}{500}
\videf{hot.tg}{3461}
\videf{sc.eden_kr}{2.3}
\videf{sc.eden_f2}{3.2}
\videf{sc.ksu_fig2}{4.2}
\videf{sc.ksu_fig1}{6.1}
\videf{sc.xu_figs5_7}{11.7}
\videf{sc.hakuta_fig3}{34.1}
\videf{sc.hakuta_fig2}{19.3}
\videf{sc.ys_ir300}{\ensuremath{-}0.5}
\videf{sc.ys_ratio_230_300}{\ensuremath{-}8.4}
\videf{sc.bs_tau_b}{\ensuremath{-}22.3}
\videf{ys.meas}{0.16}
\videf{ys.err}{0.02}
\videf{ysr.meas}{1.30}
\videf{ysr.err}{0.06}
\videf{sc.train.lo}{2.3}
\videf{sc.train.hi}{6.1}
\videf{sc.bs.meas}{9.0}
\videf{sc.bs.err}{0.6}
\videf{tah.g0.meas}{6.39}
\videf{tah.g0}{6.10}
\videf{tah.g0.x}{0.95}
\videf{tah.isat.meas}{1.4}
\videf{tah.isat}{1.54}
\videf{tah.isat.err}{0.1}
\videf{tah.isat.x}{1.10}
\videf{tah.eta.meas}{0.27}
\videf{tah.eta}{0.285}
\videf{tah.eta.x}{1.06}
\videf{tah.form}{0.318}
\videf{tah.tg}{448}
\videf{tah.pin}{330}
\videf{tah.torr}{665}
\videf{tah.f2torr}{2.8}
\videf{tah.groups}{kr2}
\videf{tah.def.g0}{5.94}
\videf{tah.def.x}{0.93}
\videf{tah.kr2noF.g0}{5.69}
\videf{tah.kr2noF.x}{0.89}
\videf{kr2.off}{4.73}
\videf{kr2.on}{9.17}
\videf{kr2.noF}{4.65}
\videf{kr2.gain}{4.44}
\videf{kr2.after}{3.56}
\videf{kr2.afterpc}{80}
\videf{kr2.burn}{665}
\videf{kr2.F}{$2.11{\times}10^{17}$}
\videf{kr2.k}{$3.0{\times}10^{-10}$}
\videf{kr2.ratio}{17}
\videf{kr2.lo.k}{$1.0{\times}10^{-11}$}
\videf{kr2.lo}{5.85}
\videf{kr2.hi.k}{$1.0{\times}10^{-9}$}
\videf{kr2.hi}{9.49}
\videf{tah.tri.g0}{3.02}
\videf{tah.tri.x}{0.47}
\videf{tri.lo.pow}{0.5}
\videf{tri.klimek.hi}{425}
\videf{tah.own.g0}{6.56}
\videf{tah.own.te}{3.8}
\videf{tah.own.ne}{$2.55{\times}10^{14}$}
\videf{tah.own.tg}{428}
\videf{tah.beam.g0}{6.10}
\videf{tah.beam.te}{1.6}
\videf{tah.beam.ne}{$6.3{\times}10^{14}$}
\videf{tah.beam.tg}{448}
\videf{tah.own.x}{1.03}
\videf{pure.Kr.x}{1.02}
\videf{pure.Ar.x}{1.03}
\videf{pure.dr.lo}{1.00}
\videf{pure.dr.hi}{1.00}
\videf{pure.dr.te.lo}{0.026}
\videf{pure.dr.te.hi}{2}
\videf{pure.k15}{$2.4{\times}10^{-31}$}
\videf{pure.conv.lo}{$2.74{\times}10^{-31}$}
\videf{pure.conv.hi}{$3.3{\times}10^{-31}$}
\videf{pure.vuv}{0.57}
\videf{pure.vuv.meas}{0.42}
\videf{pure.vuv.err}{0.05}
\videf{pure.vuv.x}{1.36}
\videf{pure.vuv.ideal}{0.54}
\videf{pure.vuv.nothermal}{0.55}
\videf{pure.thermal}{0.060}

% LIGHT v0_1 accuracy numbers (slow). GENERATED by make_overview_v0_1.py
% from LIGHT v0_1 (0.1.1) and ML/. Do not edit. Rerun the script instead.
\providecommand\videf[2]{\expandafter\def\csname vi@#1\endcsname{#2}}
\videf{fit.k_BC.lo}{$4.63{\times}10^{-12}$}
\videf{fit.k_BC.mid}{$9.8{\times}10^{-12}$}
\videf{fit.k_BC.hi}{$1.73{\times}10^{-11}$}
\videf{fit.K_eq.lo}{0.84}
\videf{fit.K_eq.mid}{1.07}
\videf{fit.K_eq.hi}{1.38}
\videf{fit.tau_C.lo}{41}
\videf{fit.tau_C.mid}{50}
\videf{fit.tau_C.hi}{63}
\videf{fit.tau_B.lo}{6.79}
\videf{fit.tau_B.mid}{7.02}
\videf{fit.tau_B.hi}{7.22}
\videf{fit.k_F2.lo}{$3.92{\times}10^{-10}$}
\videf{fit.k_F2.mid}{$4.57{\times}10^{-10}$}
\videf{fit.k_F2.hi}{$5.3{\times}10^{-10}$}
\videf{fit.k_F2_C.lo}{$2.74{\times}10^{-10}$}
\videf{fit.k_F2_C.mid}{$3.02{\times}10^{-10}$}
\videf{fit.k_F2_C.hi}{$3.32{\times}10^{-10}$}
\videf{fit.k3_Kr.lo}{$1.26{\times}10^{-30}$}
\videf{fit.k3_Kr.mid}{$1.92{\times}10^{-30}$}
\videf{fit.k3_Kr.hi}{$2.44{\times}10^{-30}$}
\videf{fit.k3_Kr_C.lo}{$1.93{\times}10^{-31}$}
\videf{fit.k3_Kr_C.mid}{$3.64{\times}10^{-31}$}
\videf{fit.k3_Kr_C.hi}{$5.12{\times}10^{-31}$}
\videf{fit.bshare.lo}{0.57}
\videf{fit.bshare.mid}{0.69}
\videf{fit.bshare.hi}{0.80}
\videf{fit.bshare.light}{0.68}
\videf{fit.n}{10}
\videf{post.g0.lo}{5.38}
\videf{post.g0.mid}{5.79}
\videf{post.g0.hi}{6.63}
\videf{post.n}{40}
\videf{link.base}{0.285}
\videf{link.yield.lo}{0.256}
\videf{link.yield.lo.d}{\ensuremath{-}0.029}
\videf{link.yield.hi}{0.299}
\videf{link.yield.hi.d}{+0.014}
\videf{link.yield.arf}{0.227}
\videf{link.yield.arf.d}{\ensuremath{-}0.058}
\videf{link.nascentC}{0.307}
\videf{link.nascentC.d}{+0.022}
\videf{link.ionY}{0.242}
\videf{link.ionY.d}{\ensuremath{-}0.043}
\videf{link.ionC}{0.268}
\videf{link.ionC.d}{\ensuremath{-}0.017}
\videf{link.ionRate}{0.283}
\videf{link.ionRate.d}{\ensuremath{-}0.002}
\videf{link.kr2F}{0.268}
\videf{link.kr2F.d}{\ensuremath{-}0.017}
\videf{link.kr2split.lo}{0.282}
\videf{link.kr2split.lo.d}{\ensuremath{-}0.003}
\videf{link.kr2split.hi}{0.288}
\videf{link.kr2split.hi.d}{+0.003}
\videf{link.k3}{0.291}
\videf{link.k3.d}{+0.006}
\videf{link.minor}{0.006}
\videf{el.loop.te}{1.61}
\videf{el.loop.ne}{$6.46{\times}10^{14}$}
\videf{el.loop.g0}{6.04}
\videf{el.loop.form}{0.316}
\videf{el.light.te}{1.60}
\videf{el.light.ne}{$6.3{\times}10^{14}$}
\videf{el.light.g0}{6.10}
\videf{el.light.form}{0.318}
\videf{el.dev.te}{0.7}
\videf{el.dev.ne}{2.5}
\videf{el.dev.g0}{0.9}
\videf{el.dev.formation}{0.5}
\videf{el.eatt}{0.55}
\videf{el.eatt.32}{2.41}
\videf{el.rpow}{0.55}
\videf{el.rcap}{0.95}
\videf{sw.krw.lo}{0.96}
\videf{sw.krw.hi}{1.03}
\videf{sw.f2.lo}{0.48}
\videf{sw.f2.hi}{2.41}
\def\viLinks{KrF$^*$ yield of \Krs{} + \Ft & 0.90 & 0.70 & the measured 0.90 less its stated 0.20 & 0.256 & \ensuremath{-}0.029\\
 &  & 1.0 & the published assignment & 0.299 & +0.014\\
 &  & 0.5 & as for Ar$^*$ + F$_2$ & 0.227 & \ensuremath{-}0.058\\
share of that KrF formed in C & 0.38 & 0 & if nascent C relaxes to B at 1 atm & 0.307 & +0.022\\
KrF$^*$ yield of ion-ion recombination & 1 & 0.64 & as for Cl$^-$ & 0.242 & \ensuremath{-}0.043\\
share of that KrF formed in C & 0 & 0.38 & as for SF$_6^-$ & 0.268 & \ensuremath{-}0.017\\
ion-ion recombination rate & Flannery & \x{0.38} & Avco's model input & 0.283 & \ensuremath{-}0.002\\
\Krts{} + F \ra KrF(B) + Kr & $3.0{\times}10^{-10}$ & 0 & no measurement exists & 0.268 & \ensuremath{-}0.017\\}

% LIGHT v0_1 energy ladder. GENERATED by make_energy_ladder.py
% from LIGHT v0_1 (0.1.1). Do not edit. Rerun the script instead.
\providecommand\laddef[2]{\expandafter\def\csname lad@#1\endcsname{#2}}
\laddef{version}{0.1.1}
\laddef{E.B}{5.856}
\laddef{DeB}{5.54}
\laddef{BoverKrF}{5.06}
\laddef{wl.G4}{248}
\laddef{hv.G4}{5.00}
\laddef{heat.G4}{0.06}
\laddef{heat.G3}{3.26}
\laddef{heat.K26}{5.54}
\laddef{heat.G5}{0.58}
\laddef{wl.C3}{420}
\laddef{hv.C3}{2.95}
\laddef{heat.C3}{1.53}
\laddef{gapBC}{120}
\laddef{notdrawn}{G6, VIB1, VIB2, VIB3 and VIB4}
\laddef{toe}{G2, PEN\_KR2 and SUPER}
\laddef{zero}{}
\laddef{groups}{\texttt{kr2} and \texttt{electron}}
\laddef{pale}{Kr\ensuremath{_2^*}}
\laddef{wl.V}{147}
\laddef{hv.V}{8.43}
\laddef{heat.V}{0.77}
\def\ladpicture{%
\definecolor{ladE}{HTML}{2A78D6}% electron gas
\definecolor{ladX}{HTML}{EDA100}% Kr*
\definecolor{ladB}{HTML}{1BAF7A}% KrF(B), 248 and 272 nm
\definecolor{ladC}{HTML}{008300}% KrF(C), 275 nm
\definecolor{ladT}{HTML}{4A3AA7}% Kr2F*, 420 nm
\definecolor{ladV}{HTML}{A63A83}% Kr2*, 147 nm
\definecolor{ladH}{HTML}{E34948}% heat
\definecolor{ladI}{HTML}{52514E}% ions
\definecolor{ladF}{HTML}{A8A69E}% F atoms
\definecolor{ladInk}{HTML}{0B0B0B}%
\definecolor{ladInk2}{HTML}{52514E}%
\definecolor{ladMuted}{HTML}{898781}%
\tikzset{
  lad grid/.style={draw=ladMuted!22, line width=0.3pt},
  lad axisline/.style={draw=ladMuted, line width=0.5pt},
  lad lev/.style={line width=1.4pt},
  lad pair/.style={line width=0.8pt, dash pattern=on 3pt off 2pt},
  lad tick/.style={line width=0.5pt},
  lad arr/.style={line width=0.7pt, -{Triangle[width=3.2pt,length=3.4pt]}},
  lad hv/.style={line width=1.3pt, -{Triangle[width=4.2pt,length=4.4pt]}},
  lad on/.style={},
  lad opt/.style={dash pattern=on 2.4pt off 1.6pt},
  lad zero/.style={dash pattern=on 0.6pt off 1.4pt, line cap=round},
  lad heat/.style={draw=ladH, line width=0.6pt, decorate,
    decoration={snake, amplitude=0.5mm, segment length=1.4mm,
    post length=0.5mm}},
  lad lab/.style={font=\scriptsize, text=ladInk, inner sep=1pt},
  lad lab2/.style={font=\scriptsize, text=ladInk, inner sep=1pt,
    align=right},
  lad lab2l/.style={font=\scriptsize, text=ladInk, inner sep=1pt,
    align=left},
  lad alab/.style={font=\scriptsize, text=ladInk, inner sep=0.8pt,
    fill=white, fill opacity=0.85, text opacity=1, align=center},
  lad slab/.style={font=\tiny, text=ladInk, inner sep=0.6pt,
    fill=white, fill opacity=0.85, text opacity=1, align=right},
  lad hlab/.style={font=\tiny, text=ladInk2, inner sep=0.8pt,
    fill=white, fill opacity=0.85, text opacity=1},
  lad tlab/.style={font=\tiny, text=ladInk2, inner sep=0.8pt},
  lad axis/.style={font=\scriptsize, text=ladInk2, inner sep=1pt},
  lad head/.style={font=\footnotesize\bfseries, text=ladInk2},
}
\begin{tikzpicture}
\draw[lad grid] (0.950,1.000) -- (16.200,1.000);
\draw[lad grid] (0.950,3.000) -- (16.200,3.000);
\draw[lad grid] (0.950,5.000) -- (16.200,5.000);
\draw[lad grid] (0.950,7.000) -- (16.200,7.000);
\draw[lad grid] (0.950,9.000) -- (16.200,9.000);
\draw[lad grid] (0.950,11.000) -- (16.200,11.000);
\draw[lad grid] (0.950,13.000) -- (16.200,13.000);
\draw[lad grid] (0.950,15.000) -- (16.200,15.000);
\draw[lad grid] (0.950,17.000) -- (16.200,17.000);
\draw[lad axisline] (0.950,0.000) -- (0.950,17.450);
\draw[lad lev,draw=ladInk2] (2.650,3.000) -- (6.300,3.000);
\draw[lad lev,draw=ladX] (2.650,12.915) -- (6.300,12.915);
\draw[lad lev,draw=ladI] (2.650,16.996) -- (6.300,16.996);
\draw[lad lev,draw=ladV!40] (5.000,12.200) -- (6.300,12.200);
\draw[lad lev,draw=ladI] (5.000,15.846) -- (6.300,15.846);
\draw[lad tick,draw=ladMuted] (3.550,14.300) -- (4.400,14.300);
\draw[lad tick,draw=ladMuted] (3.730,14.998) -- (4.400,14.998);
\draw[lad tick,draw=ladMuted] (3.910,15.750) -- (4.400,15.750);
\draw[lad tick,draw=ladMuted] (5.740,3.770) -- (6.020,3.770);
\draw[lad pair,draw=ladInk2] (10.200,14.396) -- (12.300,14.396);
\draw[lad pair,draw=ladInk2] (10.200,13.246) -- (12.300,13.246);
\draw[lad pair,draw=ladInk2] (7.000,3.800) -- (12.300,3.800);
\draw[lad lev,draw=ladB] (8.400,8.886) -- (11.050,8.886);
\draw[lad lev,draw=ladC] (8.400,8.811) -- (11.050,8.811);
\draw[lad lev,draw=ladT] (7.300,8.276) -- (8.650,8.276);
\draw[lad tick,draw=ladMuted] (9.210,3.856) -- (9.490,3.856);
\draw[lad tick,draw=ladMuted] (9.610,4.296) -- (9.890,4.296);
\draw[lad tick,draw=ladMuted] (10.010,4.331) -- (10.290,4.331);
\draw[lad tick,draw=ladMuted] (7.840,5.326) -- (8.120,5.326);
\draw[lad lev,draw=ladInk2] (13.900,3.000) -- (15.100,3.000);
\draw[lad lev,draw=ladF] (13.900,3.800) -- (15.100,3.800);
\draw[lad pair,draw=ladInk2] (13.900,4.600) -- (15.100,4.600);
\draw[lad pair,draw=ladInk2] (13.900,1.200) -- (15.100,1.200);
\draw[lad lev,draw=ladI] (13.900,0.400) -- (15.100,0.400);
\draw[lad tick,draw=ladMuted] (14.290,6.160) -- (15.100,6.160);
\draw[lad tick,draw=ladMuted] (14.570,7.340) -- (15.100,7.340);
\draw[lad arr,lad on,draw=ladE] (2.800,3.000) -- node[pos=0.52,rotate=90,lad alab] {C2} (2.800,16.996);
\draw[lad arr,lad on,draw=ladE] (3.080,3.000) -- node[pos=0.40,rotate=90,lad alab] {G1} (3.080,12.915);
\draw[lad arr,lad opt,draw=ladE] (3.360,12.915) -- node[pos=0.78,rotate=90,lad alab] {SUPER} (3.360,3.000);
\draw[lad arr,lad on,draw=ladX] (3.080,12.915) -- node[pos=0.36,rotate=90,lad alab] {G2} (3.080,16.996);
\draw[lad arr,lad opt,draw=ladE] (3.360,12.915) -- node[pos=0.62,rotate=90,lad alab] {KSTEP} (3.360,16.996);
\draw[lad arr,lad on,draw=ladE] (3.620,3.000) -- (3.620,14.300);
\draw[lad arr,lad on,draw=ladE] (3.800,3.000) -- node[pos=0.55,rotate=90,lad alab] {C1a--c} (3.800,14.998);
\draw[lad arr,lad on,draw=ladE] (3.980,3.000) -- (3.980,15.750);
\draw[lad heat] (4.400,15.750) -- (4.400,12.915);
\draw[lad arr,lad on,draw=ladH] (2.700,12.915) -- (2.350,12.065);
\draw[lad arr,lad opt,draw=ladX] (5.100,12.915) -- (5.100,12.200);
\draw[lad arr,lad on,draw=ladI] (5.250,16.996) -- node[pos=0.50,rotate=90,lad alab] {K15} (5.250,15.846);
\draw[lad arr,lad on,draw=ladI] (5.800,15.846) -- node[pos=0.45,rotate=90,lad alab] {R11p} (5.800,12.915);
\draw[lad arr,lad opt,draw=ladV] (6.120,12.200) -- node[pos=0.62,rotate=90,lad alab] {PEN\_KR2} (6.120,15.846);
\draw[lad hv,lad opt,draw=ladV] (5.880,12.200) -- node[pos=0.50,rotate=90,lad alab] {KR2\_RAD, 147 nm} (5.880,3.770);
\draw[lad heat] (5.880,3.770) -- (5.880,3.000);
\draw[lad arr,lad opt,draw=ladV] (6.300,12.200) -- node[pos=0.42,sloped,lad alab] {K2F\_BR} (7.620,8.276);
\draw[lad arr,lad opt,draw=ladV] (6.300,12.200) -- node[pos=0.42,sloped,lad alab] {NK7} (8.660,8.856);
\draw[lad arr,lad opt,draw=ladV] (6.300,12.200) -- node[pos=0.62,sloped,lad alab] {K2F\_F} (8.900,8.856);
\draw[lad arr,lad on,draw=ladX] (6.300,12.915) -- node[pos=0.40,sloped,lad alab] {G3, G3C} (9.140,8.856);
\draw[lad arr,lad on,draw=ladI] (10.450,13.246) -- node[pos=0.50,rotate=90,lad alab] {K108} (10.450,8.856);
\draw[lad arr,lad on,draw=ladI] (10.720,14.396) -- node[pos=0.62,rotate=90,lad alab] {K26} (10.720,8.856);
\draw[lad arr,lad on,draw=ladB] (8.520,8.856) -- (8.520,8.276);
\draw[lad hv,lad on,draw=ladB] (9.350,8.856) -- node[pos=0.50,rotate=90,lad alab] {G4, 248 nm} (9.350,3.856);
\draw[lad hv,lad on,draw=ladB] (9.750,8.856) -- node[pos=0.50,rotate=90,lad alab] {G4A, 272 nm} (9.750,4.296);
\draw[lad heat] (9.750,4.296) -- (9.750,3.800);
\draw[lad hv,lad on,draw=ladC] (10.150,8.841) -- node[pos=0.50,rotate=90,lad alab] {G4C, 275 nm} (10.150,4.331);
\draw[lad heat] (10.150,4.331) -- (10.150,3.800);
\draw[lad arr,lad on,draw=ladH] (10.930,8.841) -- (10.930,8.041);
\draw[lad hv,lad on,draw=ladT] (7.980,8.276) -- node[pos=0.64,rotate=90,lad alab] {C3, 420 nm} (7.980,5.326);
\draw[lad heat] (7.980,5.326) -- (7.980,3.800);
\draw[lad arr,lad on,draw=ladH] (7.520,8.276) -- (7.520,7.476);
\draw[lad arr,lad on,draw=ladE] (14.360,3.000) -- node[pos=0.62,rotate=90,lad alab] {DIS\_a3} (14.360,6.160);
\draw[lad arr,lad on,draw=ladE] (14.640,3.000) -- node[pos=0.66,rotate=90,lad alab] {DIS\_A1} (14.640,7.340);
\draw[lad heat] (15.100,7.340) -- (15.100,4.600);
\draw[lad arr,lad on,draw=ladE] (14.080,3.000) -- node[pos=0.50,rotate=90,lad alab] {ATT} (14.080,1.200);
\draw[lad arr,lad on,draw=ladF] (14.880,4.600) -- node[pos=0.25,rotate=90,lad alab] {FF} (14.880,3.000);
\node[lad axis,anchor=east] at (0.870,1.000) {\ensuremath{-}2};
\node[lad axis,anchor=east] at (0.870,3.000) {0};
\node[lad axis,anchor=east] at (0.870,5.000) {2};
\node[lad axis,anchor=east] at (0.870,7.000) {4};
\node[lad axis,anchor=east] at (0.870,9.000) {6};
\node[lad axis,anchor=east] at (0.870,11.000) {8};
\node[lad axis,anchor=east] at (0.870,13.000) {10};
\node[lad axis,anchor=east] at (0.870,15.000) {12};
\node[lad axis,anchor=east] at (0.870,17.000) {14};
\node[lad axis,anchor=east] at (0.870,18.000) {eV};
\node[lad head,anchor=center] at (4.475,18.000) {Kr and Kr\ensuremath{_2}};
\node[lad head,anchor=center] at (9.650,18.000) {KrF and Kr\ensuremath{_2}F};
\node[lad head,anchor=center] at (14.500,18.000) {F\ensuremath{_2}};
\node[lad lab,anchor=east] at (2.570,3.000) {Kr};
\node[lad lab,anchor=east] at (2.570,12.915) {Kr\ensuremath{^*} 9.915};
\node[lad lab,anchor=east] at (2.570,16.996) {Kr\ensuremath{^+} 13.996};
\node[lad lab2l,anchor=north west] at (5.030,12.150) {Kr\ensuremath{_2^*}\\9.2};
\node[lad lab2,anchor=south east] at (6.300,15.896) {Kr\ensuremath{_2^+}\\12.846};
\node[lad tlab,anchor=west] at (4.460,14.300) {11.30};
\node[lad tlab,anchor=west] at (4.460,14.998) {11.998};
\node[lad tlab,anchor=west] at (4.460,15.750) {12.75};
\node[lad alab,anchor=east] at (2.310,12.195) {G3X};
\node[lad alab,anchor=east] at (5.040,12.557) {AK8\_Kr};
\node[lad hlab,anchor=west] at (5.980,3.370) {0.77};
\node[lad lab,anchor=west] at (12.380,14.396) {Kr\ensuremath{^+} + F\ensuremath{^-} 11.396};
\node[lad lab,anchor=west] at (12.380,13.246) {Kr\ensuremath{_2^+} + F\ensuremath{^-} 10.246};
\node[lad lab,anchor=south east] at (12.300,3.840) {Kr + F 0.8};
\node[lad lab2l,anchor=west] at (11.130,8.856) {KrF(B) 5.856\\KrF(C) 120\,cm\ensuremath{^{-1}} below\\BC, CB mix them};
\node[lad lab2,anchor=east] at (7.220,8.276) {Kr\ensuremath{_2}F\ensuremath{^*}\\5.276};
\node[lad slab,anchor=east] at (8.470,8.566) {G5,\\G5C};
\node[lad hlab,anchor=north] at (9.350,3.420) {X, 0.06};
\node[lad hlab,anchor=north] at (9.750,3.720) {A, 0.50};
\node[lad hlab,anchor=north] at (10.150,3.420) {A, 0.53};
\node[lad alab,anchor=north west] at (10.620,7.991) {QF2, QKR, QE\\QF2C, QKRC, QEC};
\node[lad hlab,anchor=north] at (7.980,3.720) {1.53};
\node[lad alab,anchor=north] at (7.520,7.426) {Q2F};
\node[lad lab,anchor=east] at (13.820,3.000) {F\ensuremath{_2}};
\node[lad lab,anchor=east] at (13.820,3.800) {F 0.8};
\node[lad lab,anchor=east] at (13.820,4.600) {F + F 1.6};
\node[lad lab,anchor=east] at (13.820,1.200) {F\ensuremath{^-} + F \ensuremath{-}1.8};
\node[lad lab,anchor=east] at (13.820,0.400) {F\ensuremath{^-} \ensuremath{-}2.6};
\node[lad tlab,anchor=west] at (15.150,6.160) {{\unsansmath\ensuremath{a\,^3\Pi_u}} 3.16};
\node[lad tlab,anchor=west] at (15.150,7.340) {{\unsansmath\ensuremath{A\,^1\Pi_u}} 4.34};
\draw[lad arr,lad on,draw=ladInk2] (1.250,-0.550) -- (1.950,-0.550);
\node[lad lab,anchor=west] at (2.050,-0.550) {a row of the default network};
\draw[lad arr,lad opt,draw=ladInk2] (8.850,-0.550) -- (9.550,-0.550);
\node[lad lab,anchor=west] at (9.650,-0.550) {a row of an optional group, off by default};
\draw[lad hv,lad on,draw=ladB] (1.250,-0.970) -- (1.950,-0.970);
\node[lad lab,anchor=west] at (2.050,-0.970) {a photon, with its row and band};
\draw[lad arr,lad on,draw=ladH] (8.850,-0.970) -- (9.550,-0.970);
\node[lad lab,anchor=west] at (9.650,-0.970) {a quench, the level's energy to heat};
\draw[lad heat] (1.250,-1.390) -- (1.950,-1.390);
\node[lad lab,anchor=west] at (2.050,-1.390) {heat as a lower state falls apart (eV) or cascades};
\draw[lad pair,draw=ladInk2] (8.850,-1.390) -- (9.550,-1.390);
\node[lad lab,anchor=west] at (9.650,-1.390) {a pair, the sum of two formation energies};
\draw[lad tick,draw=ladMuted] (1.250,-1.810) -- (1.950,-1.810);
\node[lad lab,anchor=west] at (2.050,-1.810) {a threshold, or where a photon leaves its lower state};
\draw[lad lev,draw=ladInk2!40] (8.850,-1.810) -- (9.550,-1.810);
\node[lad lab,anchor=west] at (9.650,-1.810) {a species that only an optional group makes};
\end{tikzpicture}}

\makeatletter
\newcommand\vm[1]{\ifcsname vm@#1\endcsname\csname vm@#1\endcsname
  \else\errmessage{no map number '#1'. Rerun make_system_map_v0_1.py}\fi}
\newcommand\vi[1]{\ifcsname vi@#1\endcsname\csname vi@#1\endcsname
  \else\errmessage{no overview number '#1'. Rerun make_overview_v0_1.py}\fi}
\newcommand\lad[1]{\ifcsname lad@#1\endcsname\csname lad@#1\endcsname
  \else\errmessage{no ladder number '#1'. Rerun make_energy_ladder.py}\fi}
\makeatother

% chemistry and units
\newcommand{\unit}[1]{\ensuremath{\,\mathrm{#1}}}
\newcommand{\el}{\ensuremath{\mathrm{e}^-}}
\newcommand{\Krs}{Kr\ensuremath{^*}}
\newcommand{\Krp}{Kr\ensuremath{^+}}
\newcommand{\Krtp}{Kr\ensuremath{_2^+}}
\newcommand{\Krts}{Kr\ensuremath{_2^*}}
\newcommand{\KtF}{Kr\ensuremath{_2}F\ensuremath{^*}}
\newcommand{\Ft}{F\ensuremath{_2}}
\newcommand{\Fm}{F\ensuremath{^-}}
\newcommand{\ra}{\ensuremath{\to}\ }
\newcommand{\hnu}{\ensuremath{h\nu}}
\newcommand{\Edep}{\ensuremath{E_\mathrm{dep}}}
\newcommand{\gIP}{\ensuremath{g_0I_\mathrm{sat}/P_\mathrm{in}}}
\newcommand{\x}[1]{\ensuremath{\times}#1}
% a factor of the measurement as a percentage of it (0.89 -> 89), to the
% factor's own precision: \pctof{x} for two decimals, \pctof[1]{x} for three
\newcommand{\pctof}[2][0]{\pgfmathparse{100*#2}%
  \pgfmathprintnumber[fixed,precision=#1]{\pgfmathresult}}

\begin{document}
\sansmath

{\LARGE\bfseries How accurate LIGHT v0\_1 is}\par
\vspace{4pt}
{\large The basic network of LIGHT v0\_1, version \vm{version}, on 29
September 2026, with its energy ladder, its system map and its reaction
rates}\par
\vspace{10pt}

\noindent\fbox{\parbox{\dimexpr\linewidth-2\fboxsep-2\fboxrule\relax}{%
\setlength{\parskip}{5pt}%
\textbf{Bottom line.} LIGHT v0\_1 reproduces the fitted KrF decays [1, 2] within \vi{sc.train.lo}
to \vi{sc.train.hi}\,\% rms and predicts decays that no fit used [4] within \vi{sc.xu_figs5_7}\,\% rms,
and at the Tahoma laser [5] its default network, with $n_e$ free, gives
\pctof{\vi{tah.def.x}}\,\% of the measured small-signal gain of
\vi{tah.g0.meas}\unit{\%/cm}.

With the optional \Krts{} group on as well, as ML scores it, v0\_1 gives
\pctof{\vi{tah.g0.x}}\,\% of the gain, \pctof{\vi{tah.isat.x}}\,\% of the
saturation intensity and \pctof{\vi{tah.eta.x}}\,\% of the formation
efficiency, and without the group's estimated \Krts{} + F row
\pctof{\vi{tah.kr2noF.x}}\,\% of the gain. The C/B emission ratio of Yu and Setser [3] is matched to \vi{sc.ys_ir300}\,\%. No fit used Tahoma's numbers and no rate was set by a laser
measurement, so the Tahoma scores are predictions.

Its largest misses are on data that conflict with the fitted sets,
\vi{sc.hakuta_fig2} and \vi{sc.hakuta_fig3}\,\% rms on the C-A decays of
Hakuta et al.\ [6] and \vi{sc.bs_tau_b}\,\% on the KrF(B) lifetime of
Burnham and Searles [7], and the VUV yield of pure krypton at low pump,
\pctof{\vi{pure.vuv.x}}\,\% of Morozov's measurement [8].

{\footnotesize [1] Eden, Waynant, Searles and Burnham, Appl.\ Phys.\
Lett.\ 32, 733 (1978). [2] Gadomski, Xu, Setser and Nelson, Chem.\ Phys.\
Lett.\ 189, 153 (1992). [3] Yu and Setser, J.\ Phys.\ Chem.\ 94, 2934
(1990). [4] Xu, Gadomski and Setser, J.\ Chem.\ Phys.\ 99, 2591 (1993).
[5] Spectra Technology, \emph{Experimental Support for KrF Laser
Modeling}, final report to Los Alamos, 19 July 1989, OSTI 10136715.
[6] Hakuta et al., Jpn.\ J.\ Appl.\ Phys.\ 29, 1194 (1990). [7] Burnham and Searles, J.\ Chem.\ Phys.\ 67, 5967 (1977).
[8] Morozov, Heindl, Kr\"ucken, Ulrich and Wieser, J.\ Appl.\ Phys.\ 103,
103301 (2008).\par}}}
\medskip

LIGHT v0\_1 is v0's working copy. It takes one input, a power deposited
uniformly in the electrons, and its basic network carries the charge
chemistry and fluorine recombination, \vi{rates.nBasic} reaction rows in
all. The KrF(B,C) block is v0's. Every KrF rate in it but one is a
measurement of that process or a joint fit to kinetic measurements of it,
and the exception is the KrF(B) radiative lifetime, Kannari's
\vm{tauB}\unit{ns} where Eden measured 6.8 $\pm$ 0.2\unit{ns}. The rows
v0\_1 changed carry measured values where a measurement exists. No rate
was set by a laser measurement. Scored against Tahoma, the one independent
Kr/F$_2$ laser, under the same beam drive as v0, v0\_1 gives the
small-signal gain at \x{\vi{tah.g0.x}} the measured value, the saturation
intensity at \x{\vi{tah.isat.x}} and the formation efficiency \gIP{} at
\x{\vi{tah.eta.x}}. That score includes the optional \Krts{} group, and
without its estimated \Krts{} + F row the gain is \x{\vi{tah.kr2noF.x}}.
The default network, without the group, gives \x{\vi{tah.def.x}}. The
largest open questions are the trimer rate above
\vi{tri.klimek.hi}\unit{K} and the rate of \Krts{} + F \ra KrF(B) + Kr,
and neither has been measured. Table~\ref{tab:score} lists every test. The comparison with v0 is in \texttt{docs/REVIEW/}.

\begin{table}[h]
\centering\footnotesize
\renewcommand{\arraystretch}{1.12}
\begin{tabularx}{\linewidth}{@{}>{\raggedright\arraybackslash}X
  >{\raggedright\arraybackslash}p{2.6cm}>{\raggedright\arraybackslash}p{2.5cm}
  >{\raggedright\arraybackslash}p{2.5cm}@{}}
\toprule
test & measured & v0\_1 & role of the data\\
\midrule
energy ledger closure, \vi{closure.runs} runs across the modes & & below
  \vi{closure.max} of \Edep & bookkeeping\\
KrF(B) decay against Kr and F$_2$, Eden 1978 & & rms \vi{sc.eden_kr} and
  \vi{sc.eden_f2}\,\% & train\\
KrF(B,C) slow decay, Gadomski 1992 Figs. 1 and 2 & & rms \vi{sc.ksu_fig1}
  and \vi{sc.ksu_fig2}\,\% & train\\
C/B emission ratio at 300\unit{K}, Yu and Setser 1990 & \vi{ys.meas}
  $\pm$ \vi{ys.err} & \vi{sc.ys_ir300}\,\% & train\\
KrF(B,C) decays, Xu 1993 Figs. 5 to 7 & & rms \vi{sc.xu_figs5_7}\,\% &
  held out\\
decay constants Xu 1993 printed, 50 and 100 Torr Kr & \vi{xu.50.meas}
  and \vi{xu.100.meas}\unit{s^{-1}} & \x{\vi{xu.50}} and \x{\vi{xu.100}} &
  structure check\\
C-A decays, Hakuta 1990 & & rms \vi{sc.hakuta_fig2} and
  \vi{sc.hakuta_fig3}\,\% & held out, in conflict\\
$\tau_B$, Burnham and Searles 1977 & \vi{sc.bs.meas} $\pm$
  \vi{sc.bs.err}\unit{ns} & \vi{sc.bs_tau_b}\,\% & held out, in conflict\\
C/B ratio 230\unit{K} over 300\unit{K}, Yu and Setser & \vi{ysr.meas}
  $\pm$ \vi{ysr.err} & \vi{sc.ys_ratio_230_300}\,\% & held out, partly
  behind the gap v0\_1 books\\
\midrule
\Krs{} + \Krp{} made per eV, pure Kr and Ar, Saito 2003 & & \x{\vi{pure.Kr.x}}
  and \x{\vi{pure.Ar.x}} & check\\
\Krtp{} + \el{} recombination, Shiu and Biondi 1977 & & \x{\vi{pure.dr.lo}}
  & v0\_1's input\\
\Krp{} + 2Kr \ra \Krtp{} + Kr, Smith et al. 1972 & \vi{pure.conv.lo} to
  \vi{pure.conv.hi} & \vi{pure.k15} & in v0\_1's input\\
VUV yield of pure Kr at low pump, Morozov 2008 & \vi{pure.vuv.meas}
  $\pm$ \vi{pure.vuv.err} & \vi{pure.vuv}, \x{\vi{pure.vuv.x}} & check\\
electron closure at Tahoma, kinetic EEDF loop & $T_e$ \vi{el.loop.te}\unit{eV}
  & \vi{el.light.te}\unit{eV} & reference\\
\midrule
Tahoma small-signal gain $g_0$ & \vi{tah.g0.meas}\unit{\%/cm} &
  \vi{tah.g0}, \x{\vi{tah.g0.x}} & device, score only\\
\quad the same, default network without \Krts & & \vi{tah.def.g0},
  \x{\vi{tah.def.x}} & device, score only\\
\quad the same, trimer on $T_g^{-1/2}$ & & \vi{tah.tri.g0},
  \x{\vi{tah.tri.x}} & device, score only\\
\quad the same, without \Krts{} + F \ra KrF & & \vi{tah.kr2noF.g0},
  \x{\vi{tah.kr2noF.x}} & device, score only\\
Tahoma saturation intensity $I_\mathrm{sat}$ & \vi{tah.isat.meas} $\pm$
  \vi{tah.isat.err}\unit{MW/cm^2} & \vi{tah.isat}, \x{\vi{tah.isat.x}} &
  device, score only\\
Tahoma formation efficiency \gIP & \vi{tah.eta.meas} & \vi{tah.eta},
  \x{\vi{tah.eta.x}} & device, score only\\
\bottomrule
\end{tabularx}
\caption{LIGHT v0\_1 against the data. Rates are in cm$^3$s$^{-1}$ or
cm$^6$s$^{-1}$. A percentage is v0\_1's deviation from the data, signed
for a single value and rms for a set. A \x{} factor is v0\_1 over the
measurement.}
\label{tab:score}
\end{table}

\section{How the numbers were made}

Each dataset has one role. Kinetic measurements of the channels being
fitted train them. Other kinetic measurements are held out and test the
fit. Laser performance only scores, and the fitting code refuses a dataset
of that kind. This is the repository's predictive-only policy. Measurements
score the model after the fact and never steer a number. Every number in
this note is written by \texttt{docs/make\_overview\_v0\_1.py} and
\texttt{docs/make\_system\_map\_v0\_1.py} from v0\_1 and from \texttt{ML/},
and a test fails when those files no longer match the model.

ML runs v0\_1 when the environment variable \texttt{LIGHT\_VERSION} names
it and v0 otherwise, and v0's own overview tests show that the switch
leaves every v0 number unchanged. v0\_1 takes only a power input. Tahoma
was pumped by an electron beam, so v0\_1 is scored there under the same
prototype beam drive as v0, a test fixture in \texttt{docs/} and not part
of the model. Section~\ref{sec:input} runs v0\_1's own input at Tahoma's
conditions.

\section{The KrF(B,C) block}

v0\_1 keeps v0's KrF(B,C) block. KrF has two emitting states. B lases at
248\unit{nm}, and C radiates slowly in a band at 275\unit{nm}. LIGHT treats
them as the two coupled levels of eq. 5 of Gadomski et al. (1992) and
reproduces that model's analytic decay rate to within \vi{eq5.dev} at
every pressure tested. The block has \vi{fit.n} constants. Nine come from
one joint fit to three measurement sets that agree with one another, the
KrF(B) decays of Eden et al. (1978), the Kansas State B/C decays (Gadomski
et al. 1992) and Yu and Setser's equilibrium C/B ratio (1990). The tenth,
$\tau_B$, stays at Kannari's \vm{tauB}\unit{ns}, inside the fit's interval
of \vi{fit.tau_B.lo} to \vi{fit.tau_B.hi}\unit{ns}. v0\_1 shares it
between the 248-nm B-X band and B-A emission at 272\unit{nm}, which by
Dunning and Hay's calculation carries \vm{pc.BA}\,\% of the B-X rate, so
the total radiative rate of B is unchanged. The fit's median and 68\,\%
interval for B-C mixing by Kr is $k_{BC}=$ \vi{fit.k_BC.mid}
[\vi{fit.k_BC.lo}, \vi{fit.k_BC.hi}] cm$^3$s$^{-1}$. The equilibrium ratio
is $K_\mathrm{eq}=$ \vi{fit.K_eq.mid} [\vi{fit.K_eq.lo}, \vi{fit.K_eq.hi}],
and the C lifetime is $\tau_C=$ \vi{fit.tau_C.mid} [\vi{fit.tau_C.lo},
\vi{fit.tau_C.hi}]\unit{ns}. LIGHT carries the best fit,
$k_{BC}=\vm{k.BC}$\unit{cm^3s^{-1}}, $K_\mathrm{eq}=$ \vm{Keq} and
$\tau_C=$ \vm{tauC}\unit{ns}. At 1\unit{atm} and 300\unit{K} the lasing B
state holds \vi{fit.bshare.mid} [\vi{fit.bshare.lo}, \vi{fit.bshare.hi}]
of the manifold, and LIGHT's rounded constants give \vi{fit.bshare.light}.

The fitted decay sets are matched to between \vi{sc.train.lo} and
\vi{sc.train.hi}\,\% rms. The held-out decays of Xu et al. (1993) are
matched to \vi{sc.xu_figs5_7}\,\%, and the two decay constants they print
come out at \x{\vi{xu.50}} and \x{\vi{xu.100}}. Hakuta et al. (1990)
measured faster C-A decays, missed by \vi{sc.hakuta_fig2} and
\vi{sc.hakuta_fig3}\,\% rms, and their data conflict with the three fitted
sets. Burnham and Searles (1977) measured $\tau_B=$
\vi{sc.bs.meas}\unit{ns}. Those two conflicts are v0's and remain.

The energy ledger now books Yu and Setser's measured B-C gap,
\vm{gapBC}\unit{cm^{-1}}, where v0 kept REBUILD's 80. It sets how
$K_\mathrm{eq}$ changes with temperature, and their C/B ratio between 230
and 300\unit{K} now comes out \vi{sc.ys_ratio_230_300}\,\%. That is no
longer an independent test. Their 120 $\pm$ 50\unit{cm^{-1}} is a
consensus of two routes, 110 $\pm$ 20 from the absolute ratio and 180
$\pm$ 30 from this very pair of temperatures.

Tahoma's gain is the sharpest test of the block, since $g_0$ is the
KrF(B) density times the stimulated cross section. Over \vi{post.n} draws
from the fit's posterior the gain is \vi{post.g0.mid}\unit{\%/cm}
[\vi{post.g0.lo}, \vi{post.g0.hi}] against the measured
\vi{tah.g0.meas}.

\section{Charge, recombination and fluorine}

The charge chemistry is part of v0\_1's basic network, and no electrons
are injected. $n_e$ starts from a seed of $\vm{ne}$\unit{cm^{-3}} and follows
ionization, attachment and recombination. At the reference point it runs
at $\vm{neF2}$\unit{cm^{-3}} while F$_2$ lasts and peaks at
$\vm{nePeak}$\unit{cm^{-3}} after F$_2$ burns out at \vm{F2burn}\unit{ns}.
\Krp{} turns into \Krtp{} at $\vm{k.K15}$\unit{cm^6s^{-1}}, the mean of
three measurements, and \Krtp{} recombines with electrons at Shiu and
Biondi's measured $\vm{k.R11p}\,(300\,\mathrm{K}/T_e)^{\vm{pow.R11p}}$
cm$^3$s$^{-1}$. Each recombination returns a \Krs{} and heats the gas. Of
the ions that recombine, \vm{ion.R11p}\,\% meet an electron and
\vm{ion.ii}\,\% meet F$^-$ and form KrF(B), which makes \vm{frm.ii}\,\%
of the KrF formed. Once the F$_2$ is gone, Penning ionization of \Krs{}
and recombination hand the same energy back and forth, so the flows
through \Krs{} and the ions pass 100\,\% of \Edep{} and most of that
energy ends as heat. The gas ends at \vm{TgEnd}\unit{K}.

F + F + Kr gives F atoms back to F$_2$ at $\vm{k.FF}$\unit{cm^6s^{-1}},
Cohen and Bott's evaluated Ar value, since no Kr measurement exists. At
the reference point it returns \vm{FF.in}\,\% of \Edep{} from the F atoms
to heat, far too slowly to refill the F$_2$.

\section{KrF formation}

The chain from \Krp{} and \Krs{} to KrF is v0's. \Krs{} + \Ft{} reacts at
$\vm{k.harp}$\unit{cm^3s^{-1}}, the mean of three first-hand measurements.
Velazco, Kolts and Setser (1976) measured the KrF$^*$ yield of that
reaction as \vm{pc.harpY}\,\%, and Yu and Setser (1990) put \vm{pc.fC}\,\%
of that KrF$^*$ in C. Tahoma measured \gIP{} $=$ \vi{tah.eta.meas}. ML
puts a 248-nm probe flux on v0\_1 and fits its $I_\mathrm{sat}$ at fixed
$g_0$, as Tahoma did. That gives \vi{tah.eta}. Its KrF formation rate in
5-eV quanta is \vi{tah.form} of the pump. The difference is KrF that forms
in C, or mixes into it, and decays there without meeting the laser field.

Table~\ref{tab:links} moves one link at a time to another value it could
take. The measured range of the harpoon yield, 0.70 to 1.0, moves \gIP{}
from \vi{link.yield.lo} to \vi{link.yield.hi} and brackets Tahoma's
\vi{tah.eta.meas}. Three links have no measurement at all. They are the
KrF$^*$ yield of ion-ion recombination with F$^-$, the share of that KrF
formed in C, and \Krts{} + F \ra KrF. Only analogs exist for the first
two, and both analogs lower \gIP{} toward the measurement. Removing the
third lowers it too. The predictive-only policy lets no analog set a link
and no link be chosen by its fit to Tahoma. The rest of
the \Krts{} route moves \gIP{} by \vi{link.minor} at most.

\begin{table}[H]
\centering\footnotesize
\renewcommand{\arraystretch}{1.12}
\begin{tabularx}{\linewidth}{@{}>{\raggedright\arraybackslash}X
  >{\raggedright\arraybackslash}p{1.8cm}>{\raggedright\arraybackslash}p{1.9cm}
  >{\raggedright\arraybackslash}p{3.3cm}r r@{}}
\toprule
link & v0\_1 & tried & what the value is & \gIP & change\\
\midrule
\viLinks
\bottomrule
\end{tabularx}
\caption{Each formation link moved alone from v0\_1's chain at Tahoma,
\vi{tah.f2torr}\unit{Torr} of F$_2$ in Kr at \vi{tah.torr}\unit{Torr} in
all, pumped at \vi{tah.pin}\unit{kW\,cm^{-3}}. v0\_1's chain gives
\vi{link.base}, and Tahoma measured \vi{tah.eta.meas}.}
\label{tab:links}
\end{table}

\section{The electron side}

v0\_1 integrates each krypton electron-impact rate over v0's EEDF with
two temperatures, a Maxwellian bulk at $T_e$ and a tail above
9.915\unit{eV} read from kinetic EEDFs. Each captured electron takes the
energy it really carries. For \Ft{} attachment at Tahoma that is
\vi{el.eatt}\unit{eV}, where $\tfrac32 T_e$ would book
\vi{el.eatt.32}\unit{eV}. \Krtp{} recombination goes as
$T_e^{-\vi{el.rpow}}$, so a recombining electron takes
\vi{el.rcap}\,$T_e$. A kinetic loop that iterates v0\_1 with a Boltzmann
solution, with the same recombination law as its sink, is the reference.
At Tahoma v0\_1's own closure lands within \vi{el.dev.te}\,\% of the
loop's $T_e$ and \vi{el.dev.ne}\,\% of its $n_e$, and within
\vi{el.dev.g0} and \vi{el.dev.formation}\,\% of its gain and formation.

The cross sections are v0's. With Biagi's krypton set the computed $W$ is
\vi{sw.krw.lo} to \vi{sw.krw.hi} times Nakamura's measured values. For
\Ft{} attachment the groups that measured the rate constant disagree
beyond their errors, and SIGLO gives between \vi{sw.f2.lo} and
\vi{sw.f2.hi} times their values.

On pure gases the beam drive's Spencer-Fano yields make \vi{pure.Kr.x}
times the \Krs{} and \Krp{} per eV that the measured $W$ and
$N_\mathrm{ex}/N_i$ give in Kr, and \vi{pure.Ar.x} times in Ar. v0\_1
takes its \Krtp{} recombination from Shiu and Biondi and its \Krp{}
conversion from a mean that includes Smith et al., so two of v0's
pure-gas checks are now v0\_1's inputs. The VUV yield of pure Kr at low
pump is still the one miss. v0\_1 gives \vi{pure.vuv} where Morozov et al.
(2008) measured \vi{pure.vuv.meas} $\pm$ \vi{pure.vuv.err}. Without the
thermal electrons' excitation it gives \vi{pure.vuv.nothermal}, close to
the one photon per excitation and per ion that the measured $W$ and
$N_\mathrm{ex}/N_i$ allow, \vi{pure.vuv.ideal}, so the shortfall is either
physics LIGHT lacks or a systematic error of the measurement.

\section{The power input}
\label{sec:input}

v0\_1 feeds its power to the electron gas. Run that way at Tahoma's
conditions, with no beam, it gives $g_0=$ \vi{tah.own.g0}\unit{\%/cm},
\x{\vi{tah.own.x}} the measured value. Its electrons then carry the power
at $T_e=$ \vi{tah.own.te}\unit{eV} and $n_e=$ \vi{tah.own.ne}\unit{cm^{-3}},
where the beam-sustained plasma sits at \vi{tah.beam.te}\unit{eV} and
\vi{tah.beam.ne}\unit{cm^{-3}}. The gain agrees, but the electron gas does
not resemble the beam's, so the agreement does not test the power input.
No discharge-pumped Kr/F$_2$ measurement has been scored.

\section{The energy ladder}

Figure~\ref{fig:ladder} places every species of v0\_1 at its formation
energy, the energy above ground-state Kr and \Ft{} at rest that
Appendix~A lists, and draws the rows that move energy between them. It
follows the ladder view of the website's excimer simulation, with
v0\_1's own numbers. KrF(B) lies \lad{DeB}\unit{eV} below the ion pair
\Krp{} + \Fm{}, the measured $D_e(B)$ that replaces REBUILD's estimate,
and \lad{BoverKrF}\unit{eV} above Kr + F. The \lad{wl.G4}-nm photon
carries \lad{hv.G4}\unit{eV} of that and leaves KrF(X) only
\lad{heat.G4}\unit{eV} above Kr + F. The harpoon G3 heats the gas by
\lad{heat.G3}\unit{eV} for each KrF(B) it makes, the ion-ion
recombination K26 by \lad{heat.K26}\unit{eV} and the trimer's formation by
\lad{heat.G5}\unit{eV}. \KtF{} emits \lad{hv.C3}\unit{eV} at
\lad{wl.C3}\unit{nm} and leaves \lad{heat.C3}\unit{eV} of lower-state
repulsion. The \Krts{} route is an optional group in v0\_1, drawn dashed.

\begin{figure}[p]
\centering
\ladpicture
\caption{The energy ladder of LIGHT v0\_1, version \lad{version}, in eV
above ground-state Kr and \Ft{} at rest. A bar is a species at its
formation energy and a dashed bar a pair, the sum of two. An arrow is a
row, labelled with its id and solid in the default network. The rows of
the optional groups \lad{groups} are dashed, and \lad{pale}, which only
they make, is pale. A photon's arrow drops by its quantum, and the red
squiggle below it is the heat its lower state gives the gas as it falls
apart. A red stub is a quench. The energy a row releases heats the gas,
except a photon's, which leaves, and that of \lad{toe}, which goes to the
electrons. KrF(B) and KrF(C) lie \lad{gapBC}\unit{cm^{-1}} apart and are
drawn as a pair of bars that every B and C row meets. The rows
\lad{notdrawn} move nothing between levels and are not drawn.}
\label{fig:ladder}
\end{figure}

\section{The system map}

The map on the next page draws v0\_1's basic network at the notes'
reference point, \vm{p}\unit{atm} of Kr with \vm{xF2}\,\% F$_2$,
\vm{P}\unit{kW\,cm^{-3}} for \vm{pulse}\unit{\mu s}, with $n_e$ seeded at
$\vm{ne}$\unit{cm^{-3}} and then free. Band widths are proportional to the
energy they carry, and every box conserves it. Of the deposited energy,
\vm{krf.eff}\,\% passes through KrF in 5-eV quanta. Of that KrF,
\vm{brf.G4}\,\% radiates at 248\unit{nm}, \vm{brf.tri}\,\% forms the
trimer and \vm{brf.e}\,\% is quenched by electrons, so \vm{out.248}\,\% of
the deposited energy leaves at 248\unit{nm}.

% the ladder is a float, which pdfpages would leave pending past the map
\clearpage
\includepdf[pages=1,fitpaper=true]{system_map_v0_1.pdf}

\section{Where v0\_1 is not yet reliable}

\begin{itemize}
\item The trimer above \vi{tri.klimek.hi}\unit{K}. KrF + 2Kr \ra \KtF{} takes
\vm{brf.tri}\,\% of the KrF at the reference point, where the gas ends at
\vm{TgEnd}\unit{K}. The $(T_g/300)^{\vm{tpow.G5}}$ law is Shui's
calculation, measured only up to \vi{tri.klimek.hi}\unit{K}. With the pure-Kr
$T_g^{-\vi{tri.lo.pow}}$ law, the weakest the data allow, the reference
point sends \vm{om.trimer}\,\% out at 248\unit{nm} instead of
\vm{out.248}\,\%. Tahoma's window sits at \vi{tah.tg}\unit{K}, and there
that law gives the gain at \x{\vi{tah.tri.x}} the measured value. Tahoma's
gain counts against the weak law as a score and sets nothing.
\item Hot gas. Recombination returns the ionization energy as heat. At
\vi{hot.p}\unit{atm} and \vi{hot.P}\unit{kW\,cm^{-3}} the gas reaches
\vi{hot.tg}\unit{K}, beyond every rate law's data, and \KtF{} and \Krtp{}
redissociation are not in the model. The report warns above 1000\unit{K}.
\item Ion-ion recombination with F$^-$. Its rate was measured only below
300\unit{Torr}, where Jacob et al. (1978) fitted KrF$^*$ pulse shapes and
v0\_1's form, screened as there, lies somewhat above theirs. Rokni et al.
(1977) found that \Krtp{} + F$^-$ forms KrF$^*$ as efficiently as \Krp{} +
F$^-$, but from relative intensities. The absolute KrF$^*$ yield, its B/C
split and the rate near 1\unit{atm} have no measurement.
Table~\ref{tab:links} shows what the analogs would do.
\item \Krts{} + F \ra KrF(B) + Kr. The optional \Krts{} group has six
rows. Four rest on measurements, although the even split of \Krts{} +
\Ft{} between its two products is Kannari's. \Krts{} + F is Kannari's
estimate, \vi{kr2.k}\unit{cm^3s^{-1}}, and \Krts{} + \Krts{} comes from
an unpublished report. At the reference point the group raises the 248-nm
share from \vi{kr2.off} to \vi{kr2.on}\,\%, and without \Krts{} + F it
gives \vi{kr2.noF}\,\%, so the whole gain rests on the estimate.
\vi{kr2.afterpc}\,\% of it comes after F$_2$ burns out. The
\vi{kr2.F}\unit{cm^{-3}} F atoms left then turn \Krts{} into KrF
\vi{kr2.ratio} times faster than \Krts{} radiates, and each KrF photon
gives its F atom back. At Tahoma the row moves the gain from
\vi{tah.kr2noF.g0} to \vi{tah.g0}\unit{\%/cm}. Kannari's value is no
upper bound. It is the low end of a range the report he cites gave, with
a condition, for the argon analog. The share saturates, \vi{kr2.hi}\,\% at
\vi{kr2.hi.k}\unit{cm^3s^{-1}}, and at \vi{kr2.lo.k} it is already
\vi{kr2.lo}\,\%. Until the rate is measured the group is a band from
\vi{kr2.noF} to \vi{kr2.hi}\,\%, like the trimer law.
\item F + F + Kr is held at its 300\unit{K} value, which overstates it in
hot gas.
\item Hakuta's C-A decays and Burnham and Searles' $\tau_B$ remain
unexplained conflicts.
\end{itemize}

\appendix
\section{Reaction rates}

Rates are the values at 300\unit{K}. The column $T$ gives the law v0\_1
applies to them. Electron-impact rows use $\langle\sigma v\rangle$ over
the EEDF, and Section A.3 lists their Maxwellian values. Each row's full
citation is in \texttt{v0\_1/network.py}. The row ids are the ones the code
and the system map use.

\footnotesize
\setlength{\tabcolsep}{3.5pt}
\renewcommand{\arraystretch}{1.12}

\subsection{The basic network, \vi{rates.nBasic} rows}
\begin{longtable}{@{}l>{\raggedright}p{4.1cm}>{\raggedright}p{3.1cm}
  >{\raggedright}p{1.9cm}>{\raggedright\arraybackslash}p{4.75cm}@{}}
\toprule
row & reaction & rate & $T$ & source\\
\midrule
\endhead
\bottomrule
\endfoot
\virBasic
\end{longtable}

\subsection{Optional groups on the basic network, \vi{rates.nGroups} rows}
The groups \vi{rates.groups} are off at the reference point and switch on
with \texttt{-{}-add}. The first four \Krts{} rows rest on measurements,
and the even split of \Krts{} + \Ft{} between its two products is
Kannari's. \Krts{} + F and \Krts{} + \Krts{} have no measurement among
the sources held.
\begin{longtable}{@{}l>{\raggedright}p{4.1cm}>{\raggedright}p{3.1cm}
  >{\raggedright}p{1.9cm}>{\raggedright\arraybackslash}p{4.75cm}@{}}
\toprule
row & reaction & rate & $T$ & source\\
\midrule
\endhead
\bottomrule
\endfoot
\virGroups
\end{longtable}

\subsection{Electron-impact rate coefficients}
Maxwellian $\langle\sigma v\rangle$ in cm$^3$s$^{-1}$ at the $T_e$ of each
column. The last column is the factor the two-temperature tail applies at
the reference point's mid-pulse state, $T_e=$ \vi{rates.teMid}\unit{eV}
with the tail at \vi{rates.rMid}\,$T_e$. The rows below the rule give the
energy an attached electron takes from the electron gas.
\begin{center}\scriptsize
\begin{tabular}{@{}llrrrrrr@{}}
\toprule
row & process & \multicolumn{5}{c}{$T_e$ (eV)} & tail\\
\cmidrule(lr){3-7}
 & & \virTe & factor\\
\midrule
\virXsec
\bottomrule
\end{tabular}
\end{center}

\subsection{Where v0\_1's v0 network differs}
The v0 network keeps the notes' six reactions with REBUILD's values, and
its optional groups, \vi{rates.v0groups}, are REBUILD's. These rows differ
from v0\_1's basic network or exist only in the v0 network. The rows
ending in \_C are the bc group's copies for KrF(C).
\begin{longtable}{@{}l>{\raggedright}p{4.1cm}>{\raggedright}p{3.1cm}
  >{\raggedright}p{1.9cm}>{\raggedright\arraybackslash}p{4.75cm}@{}}
\toprule
row & reaction & rate & $T$ & source\\
\midrule
\endhead
\bottomrule
\endfoot
\virVzero
\end{longtable}

\needspace{20\baselineskip}
\subsection{Formation energies}
Energies above ground-state Kr and F$_2$ at rest. Every reaction's energy
release is the difference of these numbers, so no row can make or lose
energy. $c_v$ is the heat capacity per particle in units of $k_B$.
\begin{center}
\begin{tabular}{@{}lrrl@{}}
\toprule
species & $E_\mathrm{form}$ (eV) & $c_v$ & what sets it\\
\midrule
\virSpecies
\bottomrule
\end{tabular}
\end{center}

\end{document}
